Inspired by previous projects (\citep{Amarante2019, Wijnen2014}), we present a fully integrated, and characterized system for microliter range fluid delivery targeted towards neuroscience behavior experiments. 
 
We provide an alternative to gravity-based passive systems, by designing a scalable and affordable solution that allows for dynamic control over reward delivery, without suffering from the long-term stability issues often associated with the former systems. Additionally, we also highlight the afforded control gained over flow-rate that such an non-passive system allows.

In addition to the mechanical design of the syringe pump, we also developed a custom PCB that not only controls the presented pump but can also serve as a generic controller for various types of stepper motor-based syringe pumps. This allows users to design and assemble their own pump configurations. Furthermore, by implementing the Harp protocol, the firmware affords precise timestamping and synchronization with other Harp devices, facilitating scalable and reliable experimental data collection pipelines.

\textcolor{blue}{\textbf{The syringe pump entails certain limitations that should be taken into account. Due to the physical size of the syringe pump, spatial constraints become more pronounced, particularly in experimental setups where multiple delivery systems must coexist within a limited rig. Additionally, syringe pumps tend to be more costly than gravity-based water delivery systems that rely on water valves, making them a less economical choice for large-scale or resource-sensitive implementations. While it is technically feasible to automate refilling using external valves, the process remains more complex than the straightforward refill procedures of traditional gravity-fed systems. Furthermore, the current mechanical design of the syringe pump restricts its compatibility to syringes with a maximum volume of 12 mL.}}

All the necessary design and production files for the system, including the PCB, firmware, and GUI, have been made freely available, along with detailed assembly instructions. By providing these files, our goal is not only to enable users to replicate the design but also to make modifications according to their specific needs.

To ensure seamless integration with already-existing experimental rigs (such as replacing gravity-based fluid delivery systems), we offer users with several options to interface with the system, affording a very flexible configuration and compatibility. By using pre-existent controllers available in the experimental setups, the syringe pump can be fully controlled through its low level interface, a trigger based protocol, the GUI or the Harp API. The integration of this API is already available in Bonsai, which is particularly relevant since it leverages the capabilities of the syringe pump, by offering a powerful tool for real-time data stream processing, video acquisition, close-loop tasks and real time data visualization, in a visual programming language environment. 

We also designed a novel and simple protocol which allowed us to provide an extensive characterization of trial-to-trial variability in the fluid delivery dynamics. An important stack of tests when considering the intended use of this system for rodent behavior experiments. Traditionally, fluid volume characterization relies on repeating multiple fluid deliveries and then weighing the total volume. However, this approach only provides the mean volume over multiple deliveries and fails to assess variability across single fluid delivery events. \textcolor{blue}{\textbf{As we show in our results, the volume of these events is variable. Such variability is likely the result of factors such as defective rods, material hysteresis, syringe construction and measurement methods. Given how difficult it may be to completely remove some of these sources of variability, it highlights how important it is to characterize fluid delivery systems at this resolution.}}  

\textcolor{blue}{\textbf{Importantly, we also tested the compatibility of the system with common neuroscience experimental demands by deploying it during rodent behavior, where we successfully demonstrated the controlled variation of delivered water reward amounts. While our study focuses on benchmarking the system with a water solution, other studies have since used the same system with other solutions (\textit{e.g.:} sucralose \citep{SilvaENEURO.0189-24.2024})}}. Finally, we conducted electrophysiological recording experiments, confirming that no electrical artifacts were observed.




 





